% Aikido: Vorbemerkung zur Prüfungsordnung (quellen/aikido/Prüfungsordnung.pdf, S. 2, Stand 01.01.2022), wortgetreu.
Aikido-Techniken müssen kampfwirksam sein, damit man bei ihrer Ausführung nicht einer Selbsttäuschung erliegt. Das Ziel ist zwar nicht, unbedingt gewinnen zu wollen, aber es gehört doch zur Selbstschulung, wenn man soweit als möglich nach der Wirksamkeit der Techniken strebt.

Die Bewegung muss aber natürlich bleiben, was dadurch gewährleistet wird, dass man sich vom Bauch als der Mitte des Körpers her bewegt (Hara-Prinzip), während die Gliedmaßen, besonders die Arme, locker bewegt werden. Bildlich gesprochen, soll die Bewegung wie Bambus oder wie eine Weide im Wind sein.

Die Durchführung der Aikido-Techniken sollte schließlich von Ki (= Urlebenskraft) beseelt sein. Was eigentlich in allen Budo-Arten – und sogar in allen guten Bewegungskünsten – gilt, ist im Ki-Budo, also auch im Aikido, besonders wichtig. Die weiche Bewegung muss einen Kern haben. Es geht nicht hauptsächlich um den gewollt explosionsartigen Ki-Ausbruch, sondern um das Ausströmen des Ki, was durch intensiv-hingebungsvolles Üben in uns entsteht. Schlussendlich sollen wir uns gegenüber dem Shinki (göttlichen Ki) öffnen.

Zu Anfang ist es besonders wichtig, dass man nicht auf eine Vielzahl von Techniken aus ist, sondern wenige Grundtechniken gut beherrscht. Eine Haltung, der es darum geht, sich möglichst schnell viele Techniken anzueignen und dadurch zügig zu höheren Graden kommen zu wollen, passt nicht zu den Prinzipien des Shinkiryu Aiki Budo / Aikido Shinki Rengo.

Zunächst sollen eine standfeste Grundhaltung und sichere Schritte (Sabaki) erlangt werden. Sodann soll gelernt werden, den Abstand (Maai) zum Angreifer (Seme oder Uke als Partner in der Angreiferrolle) abzuschätzen sowie auch das Timing beim Einsatz der Technik. Mit der Zeit sollen die Bewegungen immer fließender und die Details der Technik immer genauer werden. Auch die Zahl der Techniken muss kontinuierlich steigen, damit Schmalspurigkeit vermieden wird. Die Prüfungen zu höheren Graden setzen auch die Techniken der niedrigeren Grade voraus. Deswegen werden bei Prüfungen immer Stichproben von Techniken unterhalb des eigentlich zu prüfenden Grades gemacht.

Es ist darauf zu achten, dass die Technik nicht wild, grob und hektisch, sondern genau, weich und dynamisch durchgeführt wird. Mit der Zeit soll dann die wahre, beherrschte Schärfe der Technik gefördert werden. Bei der Übung sollte man immer auf den Partner / die Partnerin je nach dem Fortschrittsgrad, der körperlichen Beschaffenheit, dem Alter usw. Rücksicht nehmen.

Bei der Dan-Graduierung werden darüber hinaus die Beherrschung des Ki, die meditative innere Haltung in Schlichtheit, Bescheidenheit und Unverzagtheit des Herzens und das Engagement für Aikido im Allgemeinen und für Aikido Shinki Rengo im Besonderen berücksichtigt.

Zur Dan-Graduierung gehört auch, dass man nicht nur selber gut Techniken beherrscht, sondern auch anderen zur Entwicklung verhelfen und sie gut beurteilen kann. Dies ist eben die Voraussetzung, dass die Dan-Graduierung zugleich zur selbständigen Graduierungsbefugnis führt.
